%% notes-tex/modeling/scaling-laws/sections/6-expression.tex
%% SOURCE: tables/floors.tex, tables/budget_curve.tex and the macros of tables/expression_numbers.tex
%% are generated by experiments/041-scaling-laws/scripts/expression_record.py from the committed 019
%% result files it names (expression_ceiling_all, morphology_noise_ceiling, short_budget_spread,
%% loss_min_vs_pearson_peak, joint_checkpoint_readout, baselines_both_label, v21_readout and
%% prediction_shrinkage_probe). The parameter-count reading and the two-store comparison are quoted
%% from notes-tex/019-simb-multimodal-expression/sections/0-strand.tex and
%% notes-tex/figure-3-gate/ (branch multimodal-phenotype-retrospective) and are named as such.
%% GENERATED FILE -- do not hand-edit.
%% SOURCE: experiments/041-scaling-laws/scripts/expression_record.py, reading experiments/019-simb-multimodal/results/the files named in the script docstring
\newcommand{\budgetFirstEpochs}{350}
\newcommand{\budgetFirstMean}{0.120}
\newcommand{\budgetLastEpochs}{9,900}
\newcommand{\budgetLastMean}{0.197}
\newcommand{\budgetThousandMean}{0.161}
\newcommand{\budgetGainPastThousand}{0.036}
\newcommand{\lossLiveRuns}{27}
\newcommand{\lossMinEpochMedian}{481}
\newcommand{\lossMinEpochMax}{3,097}
\newcommand{\lossPeakEpochMedian}{2,433}
\newcommand{\lossFinalExcessPct}{9.0}
\newcommand{\lossAtPeakExcessPct}{6.8}
\newcommand{\lossGainAfterMin}{+0.057}
\newcommand{\jointExprDiff}{-0.0012}
\newcommand{\jointExprPos}{5}
\newcommand{\jointExprN}{11}
\newcommand{\jointExprRef}{0.094}
\newcommand{\jointProtDiff}{-0.0172}
\newcommand{\jointProtPos}{1}
\newcommand{\jointProtMargin}{0.015}
\newcommand{\exprBar}{0.110}
\newcommand{\exprBarSd}{0.017}
\newcommand{\exprBarSeeds}{12}
\newcommand{\vBestArm}{S\_mask}
\newcommand{\vBestMean}{0.094}
\newcommand{\vBestSd}{0.028}
\newcommand{\vRefMean}{0.063}
\newcommand{\vRefSd}{0.037}
\newcommand{\shrinkSpreadMin}{1.4}
\newcommand{\shrinkSpreadMax}{2.2}
\newcommand{\shrinkRuns}{3}
\newcommand{\shrinkProtSpreadMin}{0.5}
\newcommand{\shrinkProtSpreadMax}{0.7}
\newcommand{\shrinkProtRuns}{3}
\newcommand{\morphRealized}{13}
\newcommand{\morphObservedBest}{0.082}


\section{What the expression and proteome rounds already pin down\secstatus{todo}}
\label{sec:expression}

\begin{keybox}
\textbf{Measured.} The replicate floor of every 019 panel in Pearson units; a
budget curve over eight arms from 350 to 9{,}900 epochs; the disagreement between
the validation loss and the metric in \lossLiveRuns\ long runs; a joint-training
null on eleven partitions; the sequence-baseline bar and the best arm of the
twelve-seed round.

\textbf{Not measured.} A parameter sweep and a data sweep on the expression
store. Every 019 round varied one architectural choice at a fixed size on a fixed
store, and the sizes that differ between rounds differ with the store and the
input stack, so no exponent can be read from the record.
\end{keybox}

Experiment 019 trained the equivariant cell graph transformer on the single-deletion
transcriptome (Kemmeren 2014), the single-deletion proteome (Messner 2023) and
morphology (Ohya 2005), alone, jointly and conditioned on one another, across about
twenty rounds. Read against Sections~\ref{sec:counting} and \ref{sec:ablation}, that
record settles the floor, exhibits the compute axis on the wrong statistic, and
leaves both scaling axes unmeasured. Each piece is listed with what it is good for.

\notesec{The floor, per panel}
Table~\ref{tab:floors} is the measurement Section~\ref{sec:counting} asks for before
any sweep, done for four panels. Kemmeren's two independent cultures per deletion
give a test-retest ceiling of a mean Pearson per feature of 0.775 across strains;
Caudal's 29 re-cultured isolates 0.604; Messner's duplicated-origin strains 0.405
against 0.696 from its wild-type replicates, two estimators that disagree because 146
pairs determine neither well; morphology 0.611 from 122 wild-type replicates. These are
in Pearson, the campaign's metric. The bound on $E$ in loss units needs the per-gene
replicate variance, which Kemmeren's store carries as a standard error and a replicate
count per value and which has not been converted; for the other panels only the
pooled reliability is on file.

%% GENERATED FILE -- do not hand-edit.
%% SOURCE: experiments/041-scaling-laws/scripts/expression_record.py, reading experiments/019-simb-multimodal/results/expression_ceiling_all.json and morphology_noise_ceiling.json
\begin{table}[htbp]
\centering\footnotesize
\caption{The replicate ceiling of each 019 panel, the bound on $E$ in Pearson units. Reliability is the per-feature fraction of across-strain variance that is signal (test-retest: the correlation of two independent measurements across strains; decomposition: one minus replicate variance over total variance), averaged over features; the ceiling is the mean of its square root, the best across-strain Pearson a perfect predictor of the signal can reach against a single noisy measurement. Two estimators are given where both exist because they disagree where few pairs determine either.}
\label{tab:floors}
\begin{tabular}{llrr}
\toprule
panel & estimator & reliability & ceiling, mean Pearson per feature \\
\midrule
Kemmeren 2014, expression & test-retest, 82 paired deletions & 0.611 & 0.775 \\
 & decomposition & -- & 0.862 \\
Caudal 2024, expression & test-retest, 29 re-cultured isolates & 0.405 & 0.604 \\
 & decomposition & 0.235 & 0.459 \\
Messner 2023, proteome & test-retest, 146 duplicated-origin ORFs & 0.200 & 0.405 \\
 & decomposition, 388 wild-type replicates & 0.501 & 0.696 \\
Ohya 2005, morphology & decomposition, 122 wild-type replicates & 0.444 & 0.611 \\
\bottomrule
\end{tabular}
\end{table}


\notesec{The compute axis, in Pearson}
Table~\ref{tab:budget-curve} is a budget sweep of the kind panel~c of the figure
draws, read off the training curves of eight arms rather than run as one: the same
runs rescored on the prefix of their own history at ten budgets. The mean rises from
\budgetFirstMean\ at \budgetFirstEpochs\ epochs to \budgetLastMean\ at
\budgetLastEpochs, and flattens: the tenfold budget past 1{,}000 epochs buys
\budgetGainPastThousand. That is the shape a power law with a floor has, and it
cannot be fit as one, because the statistic is a correlation and not a loss. The
document's second control (Section~\ref{sec:ablation}) exists for exactly this case.

%% GENERATED FILE -- do not hand-edit.
%% SOURCE: experiments/041-scaling-laws/scripts/expression_record.py, reading experiments/019-simb-multimodal/results/short_budget_spread.json (W&B torchcell_019_expr_v9)
\begin{table}[htbp]
\centering\footnotesize
\caption{The compute axis as the 019 campaign measured it: the same eight long-budget arms of the v9 round rescored on the prefix of their own curves at each budget (the maximum of a centered five-epoch rolling mean of the validation Pearson per feature over epochs up to the budget). Mean and sd over the eight arms; the spread is an arm spread plus nondeterminism, so it bounds the replicate spread from above. The statistic is Pearson, not the loss, for the reason given in the text.}
\label{tab:budget-curve}
\begin{tabular}{rrrrrr}
\toprule
epochs & arms & mean Pearson per feature & sd & min & max \\
\midrule
350 & 8 & 0.1198 & 0.0096 & 0.1029 & 0.1317 \\
500 & 8 & 0.1223 & 0.0058 & 0.1130 & 0.1317 \\
700 & 8 & 0.1411 & 0.0166 & 0.1207 & 0.1606 \\
1,000 & 8 & 0.1609 & 0.0099 & 0.1485 & 0.1780 \\
1,500 & 8 & 0.1670 & 0.0117 & 0.1536 & 0.1839 \\
2,000 & 8 & 0.1745 & 0.0151 & 0.1547 & 0.1950 \\
2,800 & 8 & 0.1821 & 0.0175 & 0.1629 & 0.2085 \\
4,000 & 8 & 0.1883 & 0.0171 & 0.1663 & 0.2145 \\
6,000 & 8 & 0.1917 & 0.0177 & 0.1663 & 0.2222 \\
9,900 & 8 & 0.1965 & 0.0222 & 0.1663 & 0.2382 \\
\bottomrule
\end{tabular}
\end{table}


\notesec{Why the loss cannot be the objective as it stands}
In the \lossLiveRuns\ long runs of the same round whose loss is not the metric, the
validation loss bottoms out at a median epoch \lossMinEpochMedian\ (the latest at
\lossMinEpochMax) and the Pearson peaks at a median epoch \lossPeakEpochMedian, with
the loss already \lossAtPeakExcessPct\,\% above its minimum and
\lossGainAfterMin\ of Pearson gained between the two; the final loss sits
\lossFinalExcessPct\,\% above its minimum. The shrinkage probe on saved validation
predictions says why: the expression predictions are \shrinkSpreadMin\ to
\shrinkSpreadMax\ times too spread for their correlation on \shrinkRuns\ checkpoints
(the proteome's run the other way, \shrinkProtSpreadMin\ to \shrinkProtSpreadMax\ on
\shrinkProtRuns), and one scalar per run brings the squared error below the per-gene
mean. The loss is
a scale problem riding on a correlation that keeps improving. A scaling fit on the
raw loss would therefore read a calibration artifact as the curve. Two repairs are
available and neither has been run as a sweep: fit on a calibrated loss (the squared
error after the one-scalar shrinkage, which is the probe's statistic), or fit on the
per-gene Pearson with the understanding that it is bounded and will bend near the
floor for a reason that is not $E$.

\notesec{The bar a curve has to start above}
On the store the twelve-seed round trains on, the nearest-neighbor average over
ProtT5 embeddings reaches \exprBar\ $\pm$ \exprBarSd\ validation Pearson per feature
over \exprBarSeeds\ split seeds. The best arm of that round, \texttt{\vBestArm},
reaches \vBestMean\ $\pm$ \vBestSd\ over twelve seeds, and the reference
\vRefMean\ $\pm$ \vRefSd. Morphology realizes \morphRealized\,\% of its ceiling. A
curve fit to a model that sits at or below a trivial predictor measures the predictor,
so on this store the first point of either sweep has to clear the neighbor average
before the second point means anything. On the larger expression store the v13
reference sits within 0.01 of the same baseline (Figure~3 gate document, branch
\file{multimodal-phenotype-retrospective}).

\notesec{The multimodal data axis}
Adding the proteome's strains and labels to the trunk is a move along $D$ on a second
modality, and the joint round measured it on eleven partitions: the joint arm against
the expression-only arm \jointExprDiff\ (\jointExprPos\ of \jointExprN\ partitions
above zero, reference \jointExprRef), and against the proteome-only arm
\jointProtDiff\ (\jointProtPos\ of \jointExprN\ above zero, against a
non-inferiority margin of \jointProtMargin). Pooling a second modality did not move
the curve. Conditioning did, by $+0.03$ to $+0.05$ over the unconditioned arm
(Figure~3 gate document), and stayed below a ridge from the measured modality, which
is the linear baseline for that axis. This is the one data-axis measurement the record
holds, and it says the second modality is worth more as an input than as more
training signal for a shared trunk.

\notesec{What the record says about $N$, and why it is not a curve}
The expression strand's own summary reads 145 runs from 0.12M to 4.7M parameters
with a correlation of $-0.129$ between $\log_{10} N$ and score
(\file{notes-tex/019-simb-multimodal-expression/sections/0-strand.tex}), and the
Figure~3 gate found that no committed script reproduces that subset while all 1{,}182
scored runs give $-0.001$. Either way the runs differ in budget, input stack, head and
objective at once, so the reading is a confounded null, not $L(N)$. The one deliberate
width increase on the fast code, four layers at width 180 against six at width 90,
collapsed to the per-gene-mean predictor on 2 of 2 seeds, which is the first control
of Section~\ref{sec:ablation} failing before the sweep could start: the learning rate
that suits width 90 does not suit width 180.

\notesec{What the record does not say about $D$}
The best per-gene scores of the campaign, the v13 reference at a window mean of 0.20
and 0.17 on split seeds 0 and 1 (\file{v13\_split\_readout.json}), were trained on
about the same number of strains as the twelve-seed round (the fig3\_core expression
rows against the 1{,}103 both-label rows), with the full four-embedding input stack
and a 6{,}000-epoch budget. The gap to the small trunk's 0.06 to 0.10 is therefore
budget and input, not $D$, and the record holds no two points that differ in $D$
alone.

\notesec{The two sweeps, launched on cabbi 2026-10-09}
Both run on the fast recipe with the mask schedule on (the one twelve-seed arm with
zero collapses) at batch 128 and the unscaled rate, one run per RTX6000 card, three
split seeds, W\&B project \file{torchcell\_019\_scaling}, stages \file{scale\_d} and
\file{scale\_n} of \file{igb\_expr\_wave5.slurm} (IGB jobs 2427637 and 2427638). The
statistic is still the per-gene Pearson at the registered window, with the loss curve
logged for the calibrated fit later.

\begin{itemize}[leftmargin=1.2em,itemsep=1pt,topsep=2pt]
\item \textbf{Data.} Nested prefixes of one split-seeded permutation of the 1{,}103
both-label training rows at 1, 1/2, 1/4, 1/8 and 1/16, validation and test untouched
(\file{train\_cgt\_multitask.py}, \texttt{train\_fraction}). The budget is 1{,}200
epochs for every fraction, so the small fractions take fewer optimizer steps (one per
epoch at 1/16); the per-epoch validation curve is what the readout reads and the step
mismatch is part of the result, not hidden by it.
\item \textbf{Parameters.} Widths 45, 90, 180 and 360 at six layers (nine heads
divide all four), the learning rate scaled as $1/\text{width}$ from $3\times10^{-4}$
at 90. Hypothesis (untested): that rule keeps the wider trunks out of the collapse the
width-180 arm fell into at the width-90 rate; it is not $\mu$P, and a width whose
rate is wrong will show as a collapse rather than a point on the curve. Width 360 is
last in the array so an out-of-memory there costs the sweep nothing else.
\end{itemize}
